\chapter{Introduction}
\label{ch:intro}

\section{Motivation}
Many measurements in science do not show the quantity of interest directly. An optical tomography system records how
light is scattered by a cell, not the cell's refractive index; a fluid experiment samples the flow at a few points and
times, not the full velocity or vorticity field. Recovering the hidden quantity from such indirect, incomplete and
noisy data is an \emph{inverse problem}. Inverse problems are typically ill-posed: many different unknowns explain
the measurements equally well, so the representation of the unknown and the prior knowledge built into it decide
which solution is found.

In computer vision, 3D Gaussian Splatting (3DGS)~\citep{kerbl20233dgs} has recently changed how scenes are
represented. Instead of a pixel or voxel grid, or a neural network evaluated at every point~\citep{mildenhall2020nerf},
a scene is a set of anisotropic Gaussians with learnable positions, shapes and amplitudes. The representation is
compact, adapts its resolution to the content through densification and pruning, and is differentiable, so it can be
fitted to measurements by gradient descent. Its dynamic extension, 4D Gaussian Splatting~\citep{wu20244dgs}, keeps
one canonical set of Gaussians and learns a deformation network that moves them over time.

Gaussians are also a natural basis for many physical fields: cell nuclei are smooth blobs, and vortices in a fluid are
localised, smooth structures that are carried along by the flow. This suggests that a Gaussian representation could
act as a useful structural prior for scientific inverse problems. Whether it does, and under which conditions, has not
been studied systematically: existing Gaussian methods for scientific imaging target specific modalities such as
X-ray computed tomography~\citep{cai2024xgaussian,zha2024r2gaussian}, while the recent InverseBench
benchmark~\citep{zheng2025inversebench} evaluates diffusion-model priors, not explicit Gaussian representations.

\section{Problem statement (summary)}
This thesis studies 3D and 4D Gaussian representations as the unknown in physics-based inverse problems. Given a
known forward model $\mathcal{A}$ that maps a physical field to measurements, and measurements
$\vect{y}=\mathcal{A}(f)+\vect{n}$, we represent $f$ by Gaussians and fit their parameters so that the simulated
measurements match $\vect{y}$. We ask when this representation helps compared with a plain grid of unknowns, how it
fails, and how its motion over time should be modelled. The full problem definition and scope are given in
Chapter~\ref{ch:problem}.

\section{Contributions of Thesis I}
\begin{enumerate}
  \item \textbf{Dataset analysis.} A study of the five InverseBench problems with respect to their suitability for
  3D and 4D Gaussian representations, leading to the choice of linear inverse scattering (3D) and Navier--Stokes (4D).
  \item \textbf{A physics-driven Gaussian pipeline.} A single implementation that renders 2D/3D Gaussians on a grid
  with a local-window evaluation, couples them to the InverseBench scattering operator (verified to match the
  benchmark implementation to within $3\times10^{-7}$ relative error) and to observations of the Navier--Stokes
  solution, and trains them with densification, pruning and an optional sparsity term.
  \item \textbf{4D Gaussians with a state-space motion model.} A 4D (2D + time) Gaussian representation whose
  deformation is modelled either by an MLP, as in 4D Gaussian Splatting, or by a Mamba selective state-space
  model~\citep{gu2023mamba} that treats each Gaussian's trajectory as a sequence.
  \item \textbf{Initial results and a failure analysis.} Baseline comparisons against pixel/voxel optimisation with
  identical physics, and the identification of a ``ghost'' failure mode caused by the operator's null space, with a
  sparsity-based remedy.
  \item \textbf{Open, reproducible code.} A public repository with a Python package, experiment configurations,
  notebooks, unit tests and all results so far (Appendix~\ref{app:repro}).
\end{enumerate}

\section{Outline}
Chapter~\ref{ch:background} introduces inverse problems, the two physical models and Gaussian splatting.
Chapter~\ref{ch:literature} reviews related work. Chapter~\ref{ch:problem} defines the problem, research questions
and scope. Chapter~\ref{ch:method} describes the pipeline and Chapter~\ref{ch:formulation} its mathematical
formulation. Chapter~\ref{ch:setup} gives the experimental setup and Chapter~\ref{ch:results} the initial results.
Chapter~\ref{ch:conclusion} discusses limitations and the plan for Thesis II.
